\section{Subgrupos, Productos e Representaciones Inducidas}

\subsection*{Subgrupos Abelianos}

\hypertarget{thm39}{}
\begin{theorem}
    \textbf{3.9.} \(A\) es abeliano \( \ \leftrightharpoons \ \forall (V,\rho)\) irreducible de \(A, \ \deg\rho= 1\). \comentario{Si \((n_j)_{j\leq k}\) son los grados de las irreducibles, \(|A |  = \#\)clases \(= k = \sum^k n_j^2 \  \leftrightharpoons \ n_j = 1, \ \forall j\leq k\) } 
\end{theorem}

\begin{corollary}
    \textbf{3.9.1.} Sean \(A \leq G\) subgrupo abeliano \(: |A| = a\) e \(|G| = g\). Entonces, \(\forall (V,\rho)\) irreducible de \(G\), \(\deg \rho \leq \nicefrac{g}{a}\)\textcolor{gray}{\(\ =[G:A]\)}.
\end{corollary}

\demo{Sean \(W \leq (V, \rho|_{_A})\) subrepresentación irreducible de \(A \) (\(\dim W = 1\), \(\blacksquare\) \hyperlink{thm39}{3.9.}) e \(V':= [\rho_s(W) : s \in G] \leq V = V'\) pues es \(\rho\)-inv. Finalmente, como \(\rho_{st}(W) = \rho_s(\rho_t(W)) = \rho_s(W), \ s \in G, \ t \in A\), de los que existen a lo sumo \(\nicefrac{g}{a}\).  }

\subsection*{Producto Directo}

\begin{definition}
    Si \((V_1, \rho^1)\) e \((V_2, \rho^2)\) son representaciones de \(G_1\) e \(G_2\) de caracteres \(\chi_1\) e \(\chi_2\) tenemos representación e caracter de \(G_1 \times G_2\) en \(V_1 \otimes V_2\) dados por:   
    \[\rho^1 \otimes \rho^2 : (s_1,s_2) \mapsto \rho^1(s_1) \otimes \rho^2(s_2) \ \ \ \text{ e } \ \ \ \chi(s_1, s_2) := \chi_1(s_1)\chi_2(s_2). \] 
\end{definition}

\begin{theorem}
    \textbf{3.10.} En \(G_1 \times G_2\) tenemos (i) \ Producto tensorial de irreducibles es irreducible e (ii) \ Toda irreducible es equivalente a algún producto tensorial de irreducibles. 
\end{theorem}

\demo{
    \[\text{(i)} \ \ \frac{1}{g} \sum_{s_1, s_2} |\chi(s_1, s_2)|^2 = \frac{1}{g_1} \sum_{s_1} |\chi_1(s_1)|^2 \cdot \frac{1}{g_2} \sum_{s_2} |\chi_2(s_2)|^2 = 1, \ \blacksquare \ \text{\hyperlink{thm25}{2.5.}};\]
    Para (ii) muestra que se \((f \mid \chi_1\chi_2) = 0 \ \leftrightharpoons \ f\equiv 0\), desplegando la suma e fijando coordenadas buscando aplicar el \(\blacksquare\) \hyperlink{thm26}{2.6.}.
%    \begin{align*}
%        \text{(ii)} \ \ f\in H : (f \mid \chi_1\chi_2)  &= \sum_{s_1, s_2} f(s_1, s_2)\cdot \overline{\chi_1(s_1)}\cdot \overline{\chi_2(s_2)}\\ 
%        &= \sum_{s_i} [(f(s_i, s_j) \mid \chi_j)]\cdot \overline{\chi_i(s_i)} = \sum_{s_i} g(s_i)\cdot \overline{\chi_i(s_i)}  \\ 
%        &=0 \ \leftrightharpoons \ g \equiv 0 \equiv f, \ }.
%    \end{align*} 
}

\subsection*{Representaciones Inducidas}

\begin{definition}
    Sean \((V, \rho)\) representación de \(G \geq H\) e \(\rho|_{_H} =: \theta : H \to \GL(W)\) de \(H\), esto es \(W\leq V\) \(\theta\)-inv. Decimos que \((V,\rho)\) es \emph{inducida}  por \((W, \theta)\) si \(V =\bigoplus W_\sigma :  W_\sigma := \rho|_{\sigma} (W), \ \sigma \in \nicefrac{G}{H}\). En particular, \(\dim V = [G:H] \cdot \dim W\). 
\end{definition}

\begin{lemma}
    \textbf{3.1.} Sean \((V,\rho)\) inducida por \((W, \theta)\), \((V', \rho')\) otra representación de \(G\) e \(f\in \hom(W, V')  : f \circ \theta_t = \rho'_t \circ f, \ \forall  t \in H\). Entonces, \(\exists! F \in \hom(V, V') : F|_{_W} = f\) e \(F\circ \rho_s = \rho_s' \circ F, \ \forall s \in G\). \comentario{Haga \(F(x) :=  \rho_s'(f(\rho_s^{-1}(x))) : x \in W_\sigma\), \cite[pág. 29]{rep}  }
\end{lemma}

\begin{theorem}
    \textbf{3.11.} Sea \((W, \theta)\) representación de \(H\leq G\). Entonces, \(\exists (V,\rho) \) representación de \(G\) inducida por \((W, \theta)\), única a menos de isomorfismo. \comentario{Véase \cite[pág. 30]{rep}}
\end{theorem}

\begin{theorem}
    \textbf{3.12.} Sean \((V,\rho)\) inducida por \((W, \theta)\) de caracteres \(\chi_\rho\) e \(\chi_\theta\), \(h:= |H|\) e \(V = \bigoplus rH : r\in G \). Entonces, 
    \[\chi_\rho(u) = \sum_{\overset{r\in R}{r^{-1}ur \in H}} \chi_\theta(r^{-1}u r)= \frac{1}{h} \sum_{\overset{s\in G}{s^{-1}us \in H}} \chi_\theta(s^{-1}u s).\]
\end{theorem}

\comentario{No admite resumen, véase \cite[pág. 30]{rep}}



